\textbf{M. Shimizu} | Engenheiro de Software Júnior \hfill Outubro 2025 -- Presente \\ % [cite: 36, 38]
\vspace{-0.4cm}
\begin{itemize}[leftmargin=*, noitemsep]
    \item Liderou o projeto TPC Facelift, desenvolvendo módulos backend orientados a eventos com Java, Quarkus e Vert.x para telemetria em tempo real de ferramentas industriais na Volkswagen. % [cite: 39]
    \item Desenvolveu APIs RESTful reativas processando altos volumes de dados de forma assíncrona com Hibernate Reactive para linhas de montagem da Volkswagen, Scania e GM. % [cite: 40]
    \item Otimizou deploys através da containerização de aplicações .NET usando Docker com Windows Containers nativos, reduzindo fricções na infraestrutura dos servidores. % [cite: 41]
    \item Garantiu a resiliência do sistema através de testes automatizados com JUnit, Mockito e Gatling. % [cite: 42, 43]
\end{itemize}

\textbf{Universidade Federal do ABC (CNPq)} | Pesquisador em Machine Learning \hfill Setembro 2024 -- Presente \\ % [cite: 46, 47, 48]
\vspace{-0.4cm}
\begin{itemize}[leftmargin=*, noitemsep]
    \item Desenvolve modelos preditivos focados na autenticação biométrica via dinâmica de digitação e impressões digitais, visando atingir níveis de acurácia no estado da arte. % [cite: 32, 51, 52]
    \item Aplica redes neurais convolucionais (CNNs) e arquiteturas DeepPrint (Inception v4) para extração de características usando PyTorch e Keras. % [cite: 53, 59]
\end{itemize}

\textbf{Evo Systems} | Analista de Sistemas Trainee / Engenheiro de Software \hfill Julho 2021 -- Fevereiro 2024 \\ % [cite: 65, 67, 69, 78, 79]
\vspace{-0.4cm}
\begin{itemize}[leftmargin=*, noitemsep]
    \item Projetou e manteve APIs REST em .NET/C\#, implementando integrações com AWS SES e S3, e realizando consultas em bancos relacionais com EF Core. % [cite: 70, 71]
    \item Integrou AWS Lambda e Amazon Rekognition para autenticação de controle de acesso facial na BASF. % [cite: 77]
    \item Desenvolveu funcionalidades cross-platform em Flutter com integrações Firebase e processamento via Bluetooth em tempo real. % [cite: 72, 73, 74]
\end{itemize}
